% !TeX program = pdflatex
% Samodzielny skrypt, wersja po audycie globalnym. Wersja do książki: grafy_i_wezly_tresc.tex
\documentclass[11pt,a4paper]{report}
\usepackage[T1]{fontenc}
\usepackage[utf8]{inputenc}
\usepackage[polish]{babel}
\usepackage{lmodern}
\usepackage{microtype}
\usepackage[margin=2.45cm]{geometry}
\usepackage{amsmath,amssymb,amsthm,mathtools}
\usepackage{enumitem}
\usepackage{xcolor}
\usepackage{tikz}
\usetikzlibrary{arrows.meta,calc,positioning}
\usepackage[unicode]{hyperref}
\hypersetup{colorlinks=true,linkcolor=blue!45!black,urlcolor=blue!45!black}
\setlength{\parindent}{0pt}
\setlength{\parskip}{4pt}
\setlength{\emergencystretch}{2em}
\setlist[enumerate]{topsep=3pt,itemsep=2pt}
\numberwithin{equation}{section}
\theoremstyle{plain}
\newtheorem{twierdzenie}{Twierdzenie}[section]
\newtheorem{lemat}[twierdzenie]{Lemat}
\newtheorem{stwierdzenie}[twierdzenie]{Stwierdzenie}
\newtheorem{wniosek}[twierdzenie]{Wniosek}
\theoremstyle{definition}
\newtheorem{definicja}[twierdzenie]{Definicja}
\newtheorem{przyklad}[twierdzenie]{Przykład}
\theoremstyle{remark}
\newtheorem{uwaga}[twierdzenie]{Uwaga}
\renewcommand{\proofname}{Dowód}
\begin{document}
\hypersetup{pageanchor=false}
\begin{titlepage}
\centering
\vspace*{3.6cm}
{\Huge\bfseries Teoria grafów i węzłów\par}
\vspace{0.45cm}
{\Large w topologii rozmaitości\par}
\vspace{0.9cm}
{\Large\color{blue!45!black} Od teorii grafów do chirurgii i rozmaitości czterowymiarowych\par}
\vfill
{\large Wersja zintegrowana, wersja po audycie globalnym — listopad 2026\par}
\vspace{1cm}
{\large Samodzielny skrypt matematyczny\par}
\end{titlepage}
\pagenumbering{roman}
\tableofcontents
\clearpage
\pagenumbering{arabic}
\hypersetup{pageanchor=true}
% Wstaw przez \input w książce posiadającej pakiety i środowiska zgodne z dotychczasowym skryptem.
\definecolor{tgNavy}{RGB}{30,52,88}
\definecolor{tgBlue}{RGB}{36,100,165}
\definecolor{tgTeal}{RGB}{18,140,132}
\definecolor{tgCoral}{RGB}{211,85,71}
\definecolor{tgGold}{RGB}{192,132,26}
\definecolor{tgPurple}{RGB}{123,83,160}
\tikzset{
  tgVertex/.style={circle,draw=tgNavy!80,fill=white,line width=.8pt,
      minimum size=6mm,inner sep=0pt,font=\small},
  tgBlueVertex/.style={tgVertex,draw=tgBlue,fill=tgBlue!13},
  tgRedVertex/.style={tgVertex,draw=tgCoral,fill=tgCoral!12},
  tgTealVertex/.style={tgVertex,draw=tgTeal,fill=tgTeal!12},
  tgEdge/.style={draw=tgNavy!82,line width=1.25pt},
  tgAccent/.style={draw=tgCoral,line width=2.5pt},
  tgTree/.style={draw=tgTeal,line width=2pt},
  tgGhost/.style={draw=tgNavy!35,densely dashed,line width=1.15pt}}

\input{../rozdzialy/teoria_grafow_i_wezlow/00-jak-czytac-wspolny-skrypt.tex}
\input{../rozdzialy/teoria_grafow_i_wezlow/01-podstawy-teorii-grafow.tex}
\input{../rozdzialy/teoria_grafow_i_wezlow/02-drzewa-mosty-i-cykle-fundamentalne.tex}
\input{../rozdzialy/teoria_grafow_i_wezlow/03-przechodzenie-po-wszystkich-krawedziach-teoria-eulera.tex}
\input{../rozdzialy/teoria_grafow_i_wezlow/04-spojnosc-przekroje-i-twierdzenie-mengera.tex}
\input{../rozdzialy/teoria_grafow_i_wezlow/05-skojarzenia-systemy-reprezentantow-i-cykle-hamiltona.tex}
\input{../rozdzialy/teoria_grafow_i_wezlow/06-macierze-incydencji-cykle-i-przekroje.tex}
\input{../rozdzialy/teoria_grafow_i_wezlow/07-graf-jako-przestrzen-topologiczna.tex}
\input{../rozdzialy/teoria_grafow_i_wezlow/08-laplasjan-grafu-dyskretna-teoria-hodge-a-i-twierdzenie-kirchho.tex}
\input{../rozdzialy/teoria_grafow_i_wezlow/09-grafy-planarne-sciany-i-wzor-eulera.tex}
\input{../rozdzialy/teoria_grafow_i_wezlow/10-kryteria-planarnosci-podzialy-krawedzi-i-kolorowanie.tex}
\input{../rozdzialy/teoria_grafow_i_wezlow/11-grafy-na-powierzchniach-grafy-dualne-i-homologia-komorkowa.tex}
\input{../rozdzialy/teoria_grafow_i_wezlow/12-wielomian-tutte-a-kolorowania-i-dualnosc-algebraiczna.tex}
\input{../rozdzialy/teoria_grafow_i_wezlow/13-triangulacje-i-grafy-dualne-kompleksow-symplicjalnych.tex}
\input{../rozdzialy/teoria_grafow_i_wezlow/14-grafy-reeba-i-funkcje-morse-a-na-powierzchniach.tex}
\input{../rozdzialy/teoria_grafow_i_wezlow/15-dyskretna-teoria-morse-a-formana.tex}
\input{../rozdzialy/teoria_grafow_i_wezlow/16-od-triangulacji-do-topologii-i-geometrii-rozmaitosci.tex}
\input{../rozdzialy/teoria_grafow_i_wezlow/17-grafy-cayleya-dzialania-na-drzewach-i-teoria-bassa-serre-a.tex}
\input{../rozdzialy/teoria_grafow_i_wezlow/18-teoria-wezlow-od-diagramow-do-rozmaitosci-czterowymiarowych.tex}
\input{../rozdzialy/teoria_grafow_i_wezlow/19-grafy-taita-wielomian-tutte-a-i-macierze-goeritza.tex}
\input{../rozdzialy/teoria_grafow_i_wezlow/20-uzupelnienia-wezlow-geometria-i-chirurgia-dehna.tex}
\input{../rozdzialy/teoria_grafow_i_wezlow/21-grafy-rozkladow-jsj-i-grupy-trojrozmaitosci.tex}
\input{../rozdzialy/teoria_grafow_i_wezlow/22-wezly-i-rozmaitosci-czterowymiarowe.tex}
\input{../rozdzialy/teoria_grafow_i_wezlow/23-grafy-plumbingowe-formy-przeciecia-i-krata-e8.tex}
\input{../rozdzialy/teoria_grafow_i_wezlow/24-rozmaitosci-seiferta-i-grafy-gwiazdziste.tex}
\input{../rozdzialy/teoria_grafow_i_wezlow/25-rachunek-kirby-ego-operacje-na-ramowanych-splotach-i-grafach.tex}
\input{../rozdzialy/teoria_grafow_i_wezlow/26-wymierna-kontrakcja-od-grafow-plumbingowych-do-egzotycznych-cz.tex}
\input{../rozdzialy/teoria_grafow_i_wezlow/27-uogolniona-wymierna-kontrakcja-i-lancuchy-cpq.tex}
\input{../rozdzialy/teoria_grafow_i_wezlow/28-wprowadzenie-do-teorii-seiberga-wittena-i-zastosowan-w-wymiarz.tex}
\input{../rozdzialy/teoria_grafow_i_wezlow/29-analityczne-podstawy-teorii-seiberga-wittena-przestrzenie-modu.tex}
\input{../rozdzialy/teoria_grafow_i_wezlow/30-analiza-seiberga-wittena-od-lokalnego-przekroju-do-niezmiennik.tex}
\input{../rozdzialy/teoria_grafow_i_wezlow/31-egzotyczna-czterorozmaitosc-parka-konstrukcja-i-topologia.tex}
\input{../rozdzialy/teoria_grafow_i_wezlow/32-niezmienniki-seiberga-wittena-i-egzotycznosc-rozmaitosci-parka.tex}
\input{../rozdzialy/teoria_grafow_i_wezlow/33-chirurgia-wzdluz-wezla-i-niezmienniki-czterorozmaitosci.tex}
\input{../rozdzialy/teoria_grafow_i_wezlow/34-od-grafow-plumbingowych-do-egzotycznych-struktur-gladkich-synt.tex}
\input{../rozdzialy/teoria_grafow_i_wezlow/35-nowoczesne-niezmienniki-i-zastosowania-teorii-wezlow.tex}
\input{../rozdzialy/teoria_grafow_i_wezlow/36-literatura-i-podsumowanie.tex}
\end{document}
